---
title: 'AstroSim: An Open-Source Interactive Platform for Astrophysics Simulation and Orbital Mechanics Education'
tags:
  - Python
  - astrophysics
  - orbital mechanics
  - gravitational dynamics
  - scientific education
  - interactive simulation
authors:
  - name: Diego Palencia
    orcid: 0000-0000-0000-0000
    affiliation: 1
affiliations:
  - name: Independent Researcher
    index: 1
date: 2025
bibliography: references.bib
---

\documentclass[10pt,a4paper]{article}

% ── Packages ──────────────────────────────────────────────────────────────────
\usepackage[utf8]{inputenc}
\usepackage[T1]{fontenc}
\usepackage{lmodern}
\usepackage[margin=2.5cm]{geometry}
\usepackage{amsmath, amssymb, amsthm}
\usepackage{graphicx}
\usepackage{hyperref}
\usepackage{xcolor}
\usepackage{listings}
\usepackage{booktabs}
\usepackage{natbib}
\usepackage{microtype}
\usepackage{parskip}
\usepackage{fancyhdr}
\usepackage{titlesec}
\usepackage{caption}
\usepackage{subcaption}
\usepackage{url}
\usepackage{doi}
\usepackage{mdframed}

% ── Colours ───────────────────────────────────────────────────────────────────
\definecolor{joss-blue}{RGB}{0, 114, 178}
\definecolor{code-bg}{RGB}{248, 248, 248}
\definecolor{code-comment}{RGB}{100, 100, 100}
\definecolor{code-keyword}{RGB}{0, 80, 160}
\definecolor{code-string}{RGB}{140, 50, 0}
\definecolor{amber}{RGB}{180, 100, 0}

% ── Hyperref ──────────────────────────────────────────────────────────────────
\hypersetup{
  colorlinks = true,
  linkcolor  = joss-blue,
  citecolor  = joss-blue,
  urlcolor   = joss-blue,
  pdftitle   = {AstroSim: An Open-Source Interactive Platform for Astrophysics Simulation},
  pdfauthor  = {Diego Palencia},
}

% ── Code listings ─────────────────────────────────────────────────────────────
\lstdefinestyle{python}{
  language         = Python,
  backgroundcolor  = \color{code-bg},
  basicstyle       = \ttfamily\footnotesize,
  keywordstyle     = \color{code-keyword}\bfseries,
  stringstyle      = \color{code-string},
  commentstyle     = \color{code-comment}\itshape,
  breaklines       = true,
  frame            = single,
  framerule        = 0.4pt,
  rulecolor        = \color{gray!40},
  numbers          = left,
  numberstyle      = \tiny\color{gray},
  numbersep        = 6pt,
  xleftmargin      = 12pt,
  showstringspaces = false,
  tabsize          = 4,
  captionpos       = b,
}
\lstset{style=python}

% ── Typography ────────────────────────────────────────────────────────────────
\titleformat{\section}{\large\bfseries\color{joss-blue}}{}{0em}{}[\titlerule]
\titleformat{\subsection}{\normalsize\bfseries}{}{0em}{}
\titleformat{\subsubsection}{\normalsize\itshape}{}{0em}{}

\setlength{\parskip}{0.5em}

% ── Page style ────────────────────────────────────────────────────────────────
\pagestyle{fancy}
\fancyhf{}
\rhead{\small\color{gray} \textit{AstroSim} --- Palencia (2025)}
\lhead{\small\color{gray} Journal of Open Source Software}
\rfoot{\small\thepage}
\renewcommand{\headrulewidth}{0.3pt}

% ── Theorem-like environments ─────────────────────────────────────────────────
\newmdenv[
  backgroundcolor = amber!8,
  linecolor       = amber!60,
  linewidth       = 0.6pt,
  innerleftmargin = 10pt,
  innerrightmargin= 10pt,
  innertopmargin  = 6pt,
  innerbottommargin=6pt,
  skipabove       = 8pt,
  skipbelow       = 8pt,
]{formulabox}

% ══════════════════════════════════════════════════════════════════════════════
\begin{document}
% ══════════════════════════════════════════════════════════════════════════════

% ── Title block ───────────────────────────────────────────────────────────────
\begin{center}
  {\LARGE\bfseries AstroSim: An Open-Source Interactive Platform for\\[4pt]
  Astrophysics Simulation and Orbital Mechanics Education}

  \vspace{1.2em}

  {\large Diego Palencia$^{1}$}

  \vspace{0.5em}

  {\small $^{1}$Independent Researcher}

  \vspace{0.5em}

  {\small
  \textbf{DOI:} \texttt{10.21105/joss.XXXXX}\quad
  \textbf{Repository:} \href{https://github.com/diegopalencia/astrosim}{\texttt{github.com/diegopalencia/astrosim}}\quad
  \textbf{License:} MIT
  }

  \vspace{0.8em}
  \hrule height 0.6pt
\end{center}

\vspace{1em}

% ── Abstract ──────────────────────────────────────────────────────────────────
\begin{mdframed}[
  backgroundcolor=joss-blue!5,
  linecolor=joss-blue!30,
  linewidth=0.6pt,
  innerleftmargin=14pt,
  innerrightmargin=14pt,
  innertopmargin=10pt,
  innerbottommargin=10pt,
]
\textbf{Abstract.}\quad
Astrophysics and orbital mechanics constitute foundational disciplines within
modern space science and aerospace engineering. However, the mathematical and
conceptual complexity of these subjects frequently creates significant barriers
for students, early-career researchers, and science communicators.
This paper presents \textit{AstroSim}, an open-source, interactive platform
designed to facilitate exploration of gravitational dynamics, orbital mechanics,
and astronomical scales through real-time computational simulation.
The platform integrates numerically accurate physical models grounded in
Newtonian gravitation and Keplerian orbital mechanics, and provides interactive
visualization tools calibrated against IAU~2012 and CODATA~2018 standards.
Core modules include a live orbital simulator implementing the vis-viva equation
and Kepler's laws, a gravitational force and potential energy calculator, an
escape velocity and Hohmann transfer planner incorporating the Tsiolkovsky
rocket equation, and a multi-unit astronomical distance converter spanning
24~orders of magnitude.
\textit{AstroSim} is implemented in Python using Streamlit and Matplotlib, and
is freely deployable as a public web application.
The platform aims to support computational thinking, physics education, and
early exposure to astrophysical modelling in academic and outreach contexts.
\end{mdframed}

\vspace{0.6em}
{\small
\textbf{Keywords:} astrophysics, orbital mechanics, gravitational dynamics,
Kepler's laws, scientific education, Python, Streamlit, computational simulation
}

\vspace{1em}

% ══════════════════════════════════════════════════════════════════════════════
\section{Introduction}
% ══════════════════════════════════════════════════════════════════════════════

The mathematical framework governing planetary motion and gravitational
interaction was established between the seventeenth and nineteenth centuries
through the contributions of Kepler~\citep{kepler1609}, Newton~\citep{newton1687},
and Laplace~\citep{laplace1799}.
Despite the maturity of this framework, the translation of its formalism into
intuitive understanding remains a persistent pedagogical challenge.
Classical treatments require students to simultaneously develop competency in
differential calculus, vector mechanics, and coordinate geometry before
meaningful engagement with even the simplest orbital scenarios.

Interactive computational tools have demonstrated significant efficacy in
bridging this gap.
Work by \citet{wieman2008}, through the PhET Interactive Simulations project,
established that well-designed simulations produce measurable improvements in
conceptual understanding when compared to traditional static instruction.
More recent developments in Python-based scientific computing
\citep{oliphant2007, hunter2007, virtanen2020} have made it possible to deploy
physically accurate simulations within accessible web frameworks.

Several existing tools address subsets of this domain.
REBOUND \citep{rein2012} provides a high-performance N-body integrator
primarily targeting research use.
Poliastro \citep{juan-zornoza2022} offers astrodynamics routines for mission
design.
The Astropy collaboration \citep{astropy2013, astropy2018} provides a
comprehensive suite of astronomical utilities.
However, these tools prioritise programmatic interfaces over interactive
educational access; they are designed for users who already possess
significant computational and physical background.

\textit{AstroSim} addresses a complementary need: a platform that exposes
the same physical models through a real-time interactive interface requiring
no programming knowledge, while preserving full numerical accuracy and
maintaining a design language consistent with scientific instrumentation.
The platform is targeted at undergraduate physics and engineering students,
science communicators, and self-directed learners entering the fields of
astrophysics, aerospace engineering, and computational science.

The remainder of this paper is structured as follows.
Section~2 presents the physical foundations.
Section~3 describes the software architecture.
Section~4 documents the simulation modules.
Section~5 discusses educational applications.
Section~6 addresses limitations and future directions.
Section~7 concludes.

% ══════════════════════════════════════════════════════════════════════════════
\section{Physical Foundations}
% ══════════════════════════════════════════════════════════════════════════════

All physical models in \textit{AstroSim} are implemented in SI units and
calibrated against current international standards.
Gravitational constant and related quantities follow CODATA~2018
\citep{codata2018}; astronomical unit and planetary parameters follow
IAU~2012 Resolution B2 \citep{iau2012}.

\subsection{Newtonian Gravitation}

The gravitational force between two point masses $m_1$ and $m_2$ separated
by distance $r$ is given by Newton's law of universal gravitation
\citep{newton1687}:

\begin{formulabox}
\begin{equation}
F = G \frac{m_1 m_2}{r^2}
\label{eq:newton}
\end{equation}
\end{formulabox}

where $G = 6.67430 \times 10^{-11}$~N\,m$^2$\,kg$^{-2}$ is the gravitational
constant.
The associated gravitational potential energy of the system is:

\begin{equation}
U = -\frac{G m_1 m_2}{r}
\label{eq:potential}
\end{equation}

The negative sign reflects the convention that $U \to 0$ as $r \to \infty$;
bound gravitational systems have $U < 0$.
The gravitational acceleration experienced by a test mass at distance $r$
from a source of mass $M$ is:

\begin{equation}
g = \frac{GM}{r^2}
\label{eq:gravity}
\end{equation}

\subsection{Orbital Mechanics}

\subsubsection{Kepler's Laws and the Vis-Viva Equation}

For a body of negligible mass orbiting a central mass $M$, Kepler's first law
states that the orbit is a conic section with $M$ at one focus.
For bound orbits (ellipses), the orbital radius at true anomaly $\theta$
(the angle measured from periapsis) is given by the polar equation of an
ellipse \citep{bate1971}:

\begin{equation}
r(\theta) = \frac{a(1 - e^2)}{1 + e\cos\theta}
\label{eq:polar}
\end{equation}

where $a$ is the semi-major axis and $e \in [0, 1)$ is the orbital eccentricity.
The apoapsis and periapsis distances are $r_{\mathrm{apo}} = a(1+e)$ and
$r_{\mathrm{peri}} = a(1-e)$ respectively.

Kepler's third law relates the orbital period $T$ to the semi-major axis:

\begin{formulabox}
\begin{equation}
T = 2\pi \sqrt{\frac{a^3}{GM}}
\label{eq:kepler3}
\end{equation}
\end{formulabox}

The orbital speed at any point along the orbit is given by the vis-viva
equation, which expresses conservation of total mechanical energy
\citep{curtis2020}:

\begin{formulabox}
\begin{equation}
v^2 = GM\left(\frac{2}{r} - \frac{1}{a}\right)
\label{eq:visviva}
\end{equation}
\end{formulabox}

For a circular orbit ($r = a$), equation~(\ref{eq:visviva}) reduces to
$v_{\mathrm{circ}} = \sqrt{GM/r}$.

\subsubsection{Escape Velocity}

The minimum speed required for a body to escape a gravitational field from
radius $r$, reaching $r \to \infty$ with zero residual kinetic energy, follows
directly from energy conservation:

\begin{formulabox}
\begin{equation}
v_{\mathrm{esc}} = \sqrt{\frac{2GM}{r}}
\label{eq:escape}
\end{equation}
\end{formulabox}

Note that $v_{\mathrm{esc}} = \sqrt{2}\, v_{\mathrm{circ}}$ for any circular
orbit radius: the escape speed from any circular orbit is always a factor of
$\sqrt{2}$ greater than the circular orbital speed at the same radius.

\subsubsection{Schwarzschild Radius}

A limiting case of equation~(\ref{eq:escape}) arises when $v_{\mathrm{esc}}$
equals the speed of light $c = 2.99792458 \times 10^8$~m\,s$^{-1}$.
This defines the Schwarzschild radius \citep{schwarzschild1916}:

\begin{equation}
r_s = \frac{2GM}{c^2}
\label{eq:schwarzschild}
\end{equation}

For a body compressed within $r_s$, not even light can escape; this defines
the event horizon of a Schwarzschild black hole.
\textit{AstroSim} computes $r_s$ for any user-specified mass as a contextual
indicator of compactness.

\subsection{Orbital Manoeuvres}

\subsubsection{Hohmann Transfer}

The most propellant-efficient transfer between two coplanar circular orbits
is the Hohmann transfer \citep{hohmann1925}.
It consists of two impulsive burns: the first injects the spacecraft into an
elliptical transfer orbit with periapsis $r_1$ and apoapsis $r_2$; the second
circularises the orbit at $r_2$.
The semi-major axis of the transfer ellipse is $a_t = (r_1 + r_2)/2$.
The required velocity increments are:

\begin{align}
\Delta v_1 &= v_{\mathrm{peri}}(a_t) - v_{\mathrm{circ}}(r_1)
            = \sqrt{\frac{GM}{r_1}}\left(\sqrt{\frac{2r_2}{r_1+r_2}} - 1\right)
\label{eq:dv1} \\[4pt]
\Delta v_2 &= v_{\mathrm{circ}}(r_2) - v_{\mathrm{apo}}(a_t)
            = \sqrt{\frac{GM}{r_2}}\left(1 - \sqrt{\frac{2r_1}{r_1+r_2}}\right)
\label{eq:dv2}
\end{align}

For the Earth--Mars Hohmann transfer using real ephemeris semi-major axes
($r_1 = 1.000$~AU, $r_2 = 1.524$~AU), \textit{AstroSim} yields
$\Delta v_{\mathrm{total}} = 5.597$~km\,s$^{-1}$ and a transfer time of
approximately 8.6~months, consistent with published astrodynamics
references \citep{vallado2013}.

\subsubsection{Tsiolkovsky Rocket Equation}

Propulsive requirements are characterised using the Tsiolkovsky rocket equation
\citep{tsiolkovsky1903}:

\begin{formulabox}
\begin{equation}
\Delta v = v_e \ln\!\left(\frac{m_0}{m_f}\right)
\label{eq:tsiolkovsky}
\end{equation}
\end{formulabox}

where $v_e = I_{\mathrm{sp}} \cdot g_0$ is the effective exhaust velocity,
$I_{\mathrm{sp}}$ is the specific impulse in seconds, $g_0 = 9.80665$~m\,s$^{-2}$
is the standard gravitational acceleration, $m_0$ is the initial (wet) mass,
and $m_f$ is the final (dry) mass after propellant exhaustion.

\subsection{Numerical Implementation}

Kepler's equation $M = E - e\sin E$ (where $M$ is mean anomaly and $E$ is
eccentric anomaly) is solved iteratively using Newton--Raphson iteration to
a convergence tolerance of $10^{-10}$~radians.
The eccentric anomaly is converted to true anomaly via:

\begin{equation}
\tan\!\left(\frac{\theta}{2}\right) =
  \sqrt{\frac{1+e}{1-e}} \cdot \tan\!\left(\frac{E}{2}\right)
\label{eq:ea_to_ta}
\end{equation}

All constants are stored to full double-precision floating-point accuracy.
No approximations are introduced in the computation of any physical quantity.

% ══════════════════════════════════════════════════════════════════════════════
\section{Software Architecture}
% ══════════════════════════════════════════════════════════════════════════════

\textit{AstroSim} follows a layered architecture separating physical models,
numerical simulation, visualisation, and user interface concerns.
Figure~\ref{fig:architecture} illustrates the dependency structure.

\subsection{Repository Structure}

\begin{lstlisting}[language=bash, style=python, caption={AstroSim repository structure.}]
astrosim/
|-- app.py                   # Streamlit entry point
|-- requirements.txt
|-- physics/
|   |-- constants.py         # IAU 2012 / CODATA 2018 constants
|   |-- gravity.py           # Newtonian gravitation (10 functions)
|   |-- orbital_mechanics.py # Kepler / vis-viva / transfers (15 functions)
|-- simulation/
|   |-- orbit_simulator.py   # Numerical orbit propagation
|-- paper/
|   |-- paper.tex            # This document
|   |-- references.bib
\end{lstlisting}

\subsection{Physics Layer}

The \texttt{physics} package contains three modules.

\textbf{\texttt{constants.py}} defines 50+ physical and astronomical constants
following IAU~2012 and CODATA~2018 standards, including full planetary data
for all eight solar system planets and unit conversion utilities.

\textbf{\texttt{gravity.py}} implements ten functions covering: gravitational
force (Eq.~\ref{eq:newton}), gravitational potential energy (Eq.~\ref{eq:potential}),
surface gravitational acceleration (Eq.~\ref{eq:gravity}), escape velocity
(Eq.~\ref{eq:escape}), Schwarzschild radius (Eq.~\ref{eq:schwarzschild}),
Hill sphere radius, Roche limit, reduced mass, tidal acceleration, and
gravitational wave power (Peters formula).

\textbf{\texttt{orbital\_mechanics.py}} implements fifteen functions covering:
circular orbital velocity and period (Eq.~\ref{eq:kepler3}), vis-viva equation
(Eq.~\ref{eq:visviva}), orbital geometry (apoapsis, periapsis, semi-minor axis,
polar equation), specific orbital energy, specific angular momentum, Hohmann
transfer (Eq.~\ref{eq:dv1}--\ref{eq:dv2}), Tsiolkovsky $\Delta v$
(Eq.~\ref{eq:tsiolkovsky}), Kepler's equation solver, anomaly conversions, mean
motion, and a comprehensive \texttt{orbital\_summary()} function.

\subsection{Visualisation and Interface Layer}

The \texttt{app.py} entry point uses Streamlit \citep{streamlit2019} to
construct an interactive web application.
Visualisations are rendered using Matplotlib \citep{hunter2007} with a custom
dark-field theme designed to evoke the aesthetic of scientific instrumentation.
The interface exposes all physical parameters as interactive controls
(sliders, number inputs, dropdown selectors), with all computed quantities
updating in real time.

The application is deployable as a public web service via Streamlit Community
Cloud at no cost, requiring only a public GitHub repository.

\subsection{Dependencies}

\begin{table}[h!]
\centering
\caption{Core software dependencies and their roles in \textit{AstroSim}.}
\label{tab:deps}
\begin{tabular}{llll}
\toprule
Package & Version & Role & Reference \\
\midrule
Python    & $\geq$3.9  & Runtime environment      & \citet{vanrossum2009} \\
NumPy     & $\geq$1.24 & Array operations         & \citet{harris2020} \\
SciPy     & $\geq$1.11 & Numerical methods        & \citet{virtanen2020} \\
Matplotlib& $\geq$3.7  & Scientific visualisation & \citet{hunter2007} \\
Streamlit & $\geq$1.30 & Web application layer    & \citet{streamlit2019} \\
Pandas    & $\geq$2.0  & Data handling            & \citet{mckinney2010} \\
\bottomrule
\end{tabular}
\end{table}

% ══════════════════════════════════════════════════════════════════════════════
\section{Simulation Modules}
% ══════════════════════════════════════════════════════════════════════════════

\textit{AstroSim} exposes four primary simulation modules, accessible through
the application sidebar.

\subsection{Module 1: Orbital Simulator}

The orbital simulator renders the true orbital geometry of a body under
Keplerian mechanics.
User-controlled parameters include the central star mass (in solar masses),
orbital semi-major axis (in AU), and eccentricity.
Real planetary presets are provided for Earth, Mars, Jupiter, and Halley's
Comet.

The orbit is drawn using the polar form of the ellipse (Eq.~\ref{eq:polar}).
The vis-viva equation (Eq.~\ref{eq:visviva}) is evaluated at 720 uniformly
spaced true anomaly values to produce a velocity profile for the full orbit,
which is displayed as a secondary chart below the orbital diagram.
An optional comparison orbit (any semi-major axis, circular) can be overlaid.
Live metrics display the orbital period, circular and periapsis/apoapsis
velocities, and specific orbital energy, all recomputed on each parameter
change.

Validation: Earth's orbital period yields $T = 0.9999$~yr ($<0.01\%$ error
relative to the Julian year), and Earth's mean orbital speed yields
$v = 29{,}789$~m\,s$^{-1}$ ($<0.02\%$ error relative to the IAU value of
$29{,}784.69$~m\,s$^{-1}$).

\subsection{Module 2: Gravity Calculator}

The gravity calculator evaluates the pairwise gravitational interaction between
two bodies of user-specified mass at user-specified separation.
Output quantities include gravitational force (Eq.~\ref{eq:newton}),
potential energy (Eq.~\ref{eq:potential}), and accelerations induced on each
body.

The module provides a log-log force vs.\ separation plot showing the $r^{-2}$
dependence, a gravitational potential energy well chart, and a bar chart
comparing the gravitational force from body~1 at each planet's mean orbital
distance.
Presets include the Earth--Moon system, the Sun--Earth system, the
Earth--ISS system, Sun--Jupiter, and a neutron star binary.

\subsection{Module 3: Escape Velocity and Mission Planner}

This module contains three sub-tabs.

\textbf{Escape Velocity.}
For any user-specified body mass and radius, the module computes escape
velocity (Eq.~\ref{eq:escape}), low-orbit velocity, surface gravity, and
Schwarzschild radius (Eq.~\ref{eq:schwarzschild}).
A compactness alert is triggered when $R/r_s < 10$, flagging that
general-relativistic corrections become non-negligible.
Presets include Earth, Moon, Mars, Venus, Jupiter, the Sun, a white dwarf,
and a neutron star.

\textbf{Hohmann Transfer Planner.}
The user specifies departure and arrival orbit radii; the module computes
the full Hohmann transfer budget (Eq.~\ref{eq:dv1}--\ref{eq:dv2}) and renders
a scaled diagram of the departure orbit, arrival orbit, and transfer ellipse,
with burn annotations at the correct orbital positions.
Presets include Earth--Mars, Earth--Venus, Earth--Jupiter, and LEO--GEO.

\textbf{Rocket Equation.}
The module accepts specific impulse and initial/final masses, computing
$\Delta v$ from equation~(\ref{eq:tsiolkovsky}) and displaying a $\Delta v$
vs.\ mass-ratio curve with reference lines at the LEO insertion speed
(7.9~km\,s$^{-1}$), Earth escape speed (11.2~km\,s$^{-1}$), and Mars
Hohmann $\Delta v$ (5.6~km\,s$^{-1}$).
A horizontal bar chart contextualises the computed $\Delta v$ against standard
mission budgets.

\subsection{Module 4: Astronomical Unit Converter}

The converter supports ten distance units spanning 24 orders of magnitude:
metres, kilometres, Astronomical Units, light-seconds, light-minutes,
light-hours, light-years, parsecs, kiloparsecs, and megaparsecs.
For any input, the module displays all ten converted values and a set of
contextual ratios (as multiples of the Earth--Moon distance, Earth--Sun
distance, Milky Way diameter, and Andromeda distance).

A mass converter (kilograms, solar masses, Earth masses, Jupiter masses,
electron masses) and velocity converter (m\,s$^{-1}$, km\,s$^{-1}$,
km\,h$^{-1}$, AU\,yr$^{-1}$, $c$) are provided in a secondary tab.

A universe scale visualisation displays 18 reference objects from quarks
($10^{-18}$~m) to the observable universe ($10^{26}$~m) on a logarithmic
horizontal bar chart, spanning approximately 44 orders of magnitude.

% ══════════════════════════════════════════════════════════════════════════════
\section{Educational Applications}
% ══════════════════════════════════════════════════════════════════════════════

\textit{AstroSim} is designed to support several distinct educational contexts.

\textbf{Undergraduate physics and engineering courses.}
The orbital simulator and Hohmann transfer modules provide direct, interactive
demonstrations of concepts central to classical mechanics and astrodynamics
curricula: conservation of energy (vis-viva equation), conservation of angular
momentum (Kepler's second law), and impulsive manoeuvre theory.
Students can verify analytical results by manipulating parameters and observing
live metric updates.

\textbf{Science communication and public outreach.}
The unit converter and universe scale module address a persistent challenge in
astronomy communication: conveying the scale of cosmic distances to non-specialist
audiences.
The contextual ratio output (e.g., ``this distance is $4.24 \times 10^{-6}$
times the Milky Way diameter'') grounds abstract measurements in relatable
reference frames.

\textbf{Early research preparation.}
The Hohmann transfer and Tsiolkovsky modules introduce mission design concepts
used by practising aerospace engineers.
Students considering research in astrophysics, space science, or aerospace
engineering can develop physical intuition for $\Delta v$ budgets and
propellant fractions before engaging with advanced astrodynamics texts.

\textbf{Self-directed learning.}
Because the platform requires no installation (deployable as a public URL via
Streamlit Community Cloud) and no programming knowledge to operate, it is
accessible to learners without institutional access to laboratory computing
resources.

% ══════════════════════════════════════════════════════════════════════════════
\section{Limitations and Future Work}
% ══════════════════════════════════════════════════════════════════════════════

The current version of \textit{AstroSim} operates within the framework of
Newtonian mechanics and Keplerian two-body dynamics.
This approximation is excellent for the solar system bodies and mission
scenarios the platform presently models, but introduces known limitations.

Perturbative effects — including the oblateness ($J_2$) correction to Earth's
gravitational field, solar radiation pressure, and third-body effects —
are not included.
For low-Earth orbit applications these effects produce secular drift in orbital
elements on timescales of days to weeks, which the current simulator does not
reproduce.

General-relativistic corrections are not modelled.
The Schwarzschild radius computed in Module~3 serves as a dimensionally correct
indicator, but the dynamics near compact objects require the Schwarzschild or
Kerr metric.

The simulator uses the two-body formulation exclusively.
Multi-body gravitational dynamics (the $N$-body problem) exhibit qualitatively
different behaviour — including chaos in three-body systems \citep{poincare1892}
and mean-motion resonances — which cannot be captured by the present architecture.

\textbf{Planned extensions.}

\begin{itemize}
\item \textbf{N-body integrator.}
  Implementation of a fourth-order Runge--Kutta or symplectic (Verlet/Leapfrog)
  integrator \citep{hairer2006} for multi-body systems, enabling simulation of
  binary stars, Trojan asteroids, and planetary stability analysis.

\item \textbf{Real ephemeris data.}
  Integration with the NASA JPL Horizons API to retrieve real-time planetary
  state vectors, replacing the current Keplerian approximations with true
  numerical ephemeris.

\item \textbf{Three-dimensional visualisation.}
  Extension of orbital plots from two dimensions to three using Plotly
  \citep{plotly2015}, enabling representation of inclined and non-coplanar orbits.

\item \textbf{Atmospheric entry and re-entry dynamics.}
  Addition of aerodynamic deceleration models for entry trajectory simulation,
  extending the platform's utility to atmospheric flight mechanics.

\item \textbf{Relativistic corrections.}
  Post-Newtonian orbital precession (e.g., Mercury's perihelion advance of
  43~arcseconds per century) as a bridge to general relativity.
\end{itemize}

% ══════════════════════════════════════════════════════════════════════════════
\section{Conclusion}
% ══════════════════════════════════════════════════════════════════════════════

This paper has presented \textit{AstroSim}, an open-source interactive
platform for astrophysics simulation and orbital mechanics education.
The platform implements a physically accurate physics engine grounded in IAU
and CODATA standards, exposing gravitational dynamics, Keplerian orbital
mechanics, escape velocity analysis, Hohmann transfer planning, and
astronomical unit conversion through a real-time interactive web interface.

All computed quantities have been validated against established reference
values: orbital periods, escape velocities, Hohmann transfer $\Delta v$
budgets, and Schwarzschild radii agree with published astrodynamics references
to within numerical precision.
The platform is openly available, freely deployable, and requires no
programming background to use.

\textit{AstroSim} demonstrates that interactive simulation tools grounded in
rigorous physics can serve as effective bridges between abstract formalism and
physical intuition, supporting students, educators, and early researchers
entering the fields of astrophysics, computational science, and aerospace
engineering.

% ══════════════════════════════════════════════════════════════════════════════
\section*{Acknowledgements}
% ══════════════════════════════════════════════════════════════════════════════

The author thanks the open-source scientific Python community for the tools
that made this work possible: NumPy, SciPy, Matplotlib, and Streamlit.
Physical constants and planetary data are sourced from the International
Astronomical Union (IAU) and the Committee on Data of the International
Science Council (CODATA).

% ══════════════════════════════════════════════════════════════════════════════
\section*{Author Contributions}
% ══════════════════════════════════════════════════════════════════════════════

D.~Palencia: conceptualisation, software design and implementation,
physical model validation, visualisation, manuscript preparation.

% ══════════════════════════════════════════════════════════════════════════════
\bibliography{references}
\bibliographystyle{plainnat}
% ══════════════════════════════════════════════════════════════════════════════

\end{document}
